\documentclass[border=10pt]{standalone}
\usepackage{tikz}
\usetikzlibrary{shapes,arrows,positioning,fit,backgrounds}


\begin{tikzpicture}[
    node distance=0.3cm,
    layer/.style={rectangle, draw, thick, fill=blue!20, minimum width=10cm, minimum height=1.2cm, align=center, font=\bfseries},
    sublayer/.style={rectangle, draw, fill=blue!10, minimum width=4.5cm, minimum height=0.8cm, align=left, font=\small},
    metadata/.style={rectangle, draw, fill=yellow!30, minimum width=4cm, align=left, font=\footnotesize},
    arrow/.style={->, thick},
]

% Data Layers at bottom
\node[sublayer] (dataA) {DATA LAYER A};
\node[sublayer, right=0.5cm of dataA] (dataB) {DATA LAYER B etc....};

% Forcings, Reanalysis, Obs below data layers
\node[below=0.3cm of dataA.south, anchor=north, font=\small] (forcings) {FORCINGS};
\node[right=0.3cm of forcings] (reanalysis) {REANALYSIS};
\node[right=0.3cm of reanalysis] (obs) {OBS};

% Estimation Layer
\node[layer, above=1cm of dataA.north west, anchor=south west] (estimation) {
    ESTIMATION LAYER:\\
    \normalfont\small -comparison\\
    -filtering\\
    -bayesian interference
};

% Cartographic Layer
\node[layer, above=of estimation] (cartographic) {CARTOGRAPHIC LAYER};

% Analysis Layer
\node[layer, above=of cartographic] (analysis) {
    ANALYSIS LAYER\\
    \normalfont\small - Operator (Time and space)\\
    -Extractor (Mask - Time and space)
};

% Model Layers
\node[sublayer, above=0.8cm of analysis.north west, anchor=south west] (modelA) {MODEL LAYER A};
\node[sublayer, right=0.5cm of modelA] (modelB) {MODEL LAYER B etc....};

% Title at top
\node[above=0.5cm of modelA.north west, anchor=south west, font=\Large\bfseries] (title) {MONOLITHIC HYDROTWIN};

% Right side: State Variable box
\node[metadata, right=1.5cm of estimation.east, anchor=west, minimum height=3cm] (statevar) {
    \textbf{EACH STATE VARIABLE is an}\\
    \textbf{ontological object}\\[0.2cm]
    ex. 3D matrix\\
    [ncells][nparams][timestep]\\[0.2cm]
    + METADATA
};

% Metadata details box
\node[metadata, below=0.3cm of statevar.south west, anchor=north west, minimum height=3.5cm] (metadata) {
    \textbf{METADATA:}\\
    - provenance (model or\\
    \hspace{0.3cm}filtering et..)\\
    - spatial scale\\
    - functional role\\
    - domain of validity\\
    \hspace{0.3cm}(time and space)\\
    - underlying physical or\\
    \hspace{0.3cm}modelling assumption
};

% Git registry box
\node[metadata, below=0.3cm of metadata.south west, anchor=north west] (git) {
    \textbf{GIT synchronized registry}\\
    - Identity\\
    - Versioning\\
    - Provenance
};

% Arrows from data to estimation
\draw[arrow] (dataA.north) -- (estimation.south -| dataA.north);
\draw[arrow] (dataB.north) -- (estimation.south -| dataB.north);

% Arrows from forcings/reanalysis/obs to data layers
\draw[arrow] (forcings.north) -- (dataA.south -| forcings.north);
\draw[arrow] (reanalysis.north) -- (dataA.south -| reanalysis.north);
\draw[arrow] (obs.north) -- (dataB.south -| obs.north);

% Arrow from estimation to cartographic
\draw[arrow] (estimation.north) -- (cartographic.south);

% Arrow from cartographic to analysis
\draw[arrow] (cartographic.north) -- (analysis.south);

% Arrows from model layers to analysis
\draw[arrow] (modelA.south) -- (analysis.north -| modelA.south);
\draw[arrow] (modelB.south) -- (analysis.north -| modelB.south);

% Arrows connecting to state variable box
\draw[arrow] (analysis.east) -- (statevar.west |- analysis.east);
\draw[arrow] (cartographic.east) -- (statevar.west |- cartographic.east);

\end{tikzpicture}
